\chapter{Electrons in crystals}
\setwlchapter{crystals}
\section{From the gas to a lattice}
For the three-dimensional unpolarized electron gas let
$4\pi(r_s a_0)^3/3=1/n$, so that $k_F=(3\pi^2 n)^{1/3}$. The Fermi
sphere gives the kinetic energy per electron $3\hbar^2 k_F^2/(10m)$.
To see the different density scaling of exchange, integrate
$4\pi e_*^2/q^2$ over two occupied spheres of equal spin, where
$e_*^2=e^2/(4\pi\epsilon_0)$. At fixed separation $q$ the overlap
volume is $\pi(4k_F+q)(2k_F-q)^2/12$ for $q\le2k_F$ and zero otherwise.
The radial measure cancels the interaction denominator. Dividing by
the density and including the exchange factor $-1/2$ gives
$E_x/N=-3e_*^2 k_F/(4\pi)$. The neutralizing background removes the
Hartree energy. Thus kinetic and exchange terms scale respectively as
$r_s^{-2}$ and $r_s^{-1}$; their competition alone does not give the
exact correlation energy.

On a Bravais lattice with hopping $-t$ on each nearest-neighbour bond,
$\epsilon_k=-2t\sum_{\alpha=1}^{d}\cos k_\alpha$ for the chain,
square and simple cubic lattices with unit spacing. Bandwidth $4dt$
does not determine the density of states: at a stationary point its
shape also depends on the signs of the local curvatures. A 2D saddle
gives the logarithmic van Hove singularity at zero energy on the square
lattice; the 1D band edge instead has an inverse-square-root divergence.
\casestudy{square}{A square-lattice Bloch matrix}

\section{Orbital basis, valleys and topology}
Take one $A$ and one $B$ orbital per cell. In the SSH chain the
off-diagonal matrix element in a cell-periodic convention is
$q(k)=v+w e^{-ik}$. Its path is a circle about $v$. With chiral
symmetry and a fixed cell and termination, this path encloses the
origin if $|w|>|v|$; one cannot change that winding without crossing
$q=0$. The spectrum $\pm|q|$ is insensitive to the orientation of
the path, but the phase of an eigenvector is not. At $v=w$ the gap
closes at $k=\pi$ for real positive hoppings.
\casestudy{ssh}{Where the dimerized-chain gap closes}
\subsection{Measuring a winding without choosing eigenvector phases}
An eigensolver may multiply its eigenvector by a different phase at
every momentum. Direct finite differences of these vectors are then
meaningless. In a two-band chiral model with no diagonal terms use
the phase of $q(k)=h_{AB}(k)$ instead. Along a closed mesh add
$\arg[q(k_{j+1})/q(k_j)]$, taking each step on the principal branch.
The unit cell and ordered pair of orbitals must stay fixed. For the
orientation $q=v+w e^{-ik}$ the nontrivial integer is negative.
The occupied-band Berry phase modulo $2\pi$ is $\pi$ times this
winding modulo $2\pi$ in the cell-periodic convention.
\begin{wl}
winding[v_,w_,n_] := Module[{q = N[v + w Exp[-2 Pi I Range[0,n-1]/n]]},
  Round[Total[Arg[RotateLeft[q]/q]]/(2 Pi)]];
TableForm[Table[{v,1,winding[v,1,40],winding[v,1,80]}, {v,{1/2,2}}],
 TableHeadings -> {None,{"v","w","40 points","80 points"}}]
\end{wl}
At $v=w$ the mesh includes $q(\pi)=0$, so the winding is undefined;
do not repair division by zero by adding a small imaginary shift.
A phase step can cross the branch cut on a coarse mesh even short of
a closing. Refining the mesh or following the circle analytically
distinguishes this sampling error from a transition.

\subsection{A second translation vector}
In graphene the two independent translation vectors supply a second
phase. At a valley the three complex hopping amplitudes sum to zero,
and linearizing their sum gives a Dirac cone. A staggered sublattice
potential gives equal mass signs in the two valleys for a fixed spin;
intrinsic Kane--Mele spin-orbit coupling gives opposite signs and
reverses them for the other spin. Time reversal therefore allows
opposite spin-resolved Chern numbers with zero total Chern number in
the spin-conserving limit.
\casestudy{graphene}{Destructive interference at a Dirac point}
\casestudy{rashba}{Which spin component splits the band?}

\section{Flux per plaquette}
In Landau gauge the forward $y$ hop at coordinate $x$ carries
$e^{2\pi i\alpha x}$ for flux $\alpha$ in units of a flux quantum per
plaquette. For $\alpha=p/q$ in lowest terms a magnetic cell contains
$q$ sites. At fixed $k_y$ the diagonal entries of the Harper matrix
are $-2t\cos(k_y+2\pi p x/q)$; nearest $x$ neighbours are coupled
by $-t$ with a boundary twist $e^{iqk_x}$. Here energies are in units
of $t$; at the centre of the magnetic Brillouin zone:
\begin{wl}
harper[q_, ky_] := Table[-2 Cos[ky + 2 Pi i/q] KroneckerDelta[i, j] -
  Boole[Mod[i-j, q] == 1] - Boole[Mod[j-i, q] == 1],
  {i, 0, q-1}, {j, 0, q-1}];
harper[3, 0] // MatrixForm
\end{wl}
The oriented product of link phases around a plaquette is gauge
invariant; individual bond phases need not be. The same distinction
matters for current operators derived by differentiating the hopping
with respect to vector potential~\cite{vucicevic}.
\begin{exercise}[A magnetic cell]\label{ex:harper}
For flux $1/3$ find the three energies at $k_x=k_y=0$. Check their
sum and the sum of their squares without using the eigenvectors.
\end{exercise}
